% Generated draft: requires mathematical and source review.
\documentclass[11pt]{article}
\usepackage[margin=25mm]{geometry}
\usepackage{fontspec}
\setmainfont{Latin Modern Roman}
\usepackage{amsmath,amssymb,mathtools,unicode-math}
\setmathfont{Latin Modern Math}
\usepackage{enumitem,needspace,xurl,hyperref}
\setlength{\parindent}{0pt}
\setlength{\parskip}{5pt}
\setlength{\emergencystretch}{3em}
\newcommand{\N}{\mathbb{N}}
\newcommand{\Z}{\mathbb{Z}}
\newcommand{\Q}{\mathbb{Q}}
\newcommand{\R}{\mathbb{R}}
\newcommand{\C}{\mathbb{C}}
\newcommand{\F}{\mathbb{F}}
\newcommand{\A}{\mathbb{A}}
\newcommand{\PP}{\mathbb{P}}
\newcommand{\E}{\mathbb{E}}
\newcommand{\eps}{\varepsilon}
\newcommand{\dd}{\,\mathrm{d}}
\begin{document}
\section*{UnsolvedMath ID: 10400043}
\textbf{Source problem code:} catalogue entry
\subsection{Definitions and mathematical statement}
The source, attributed to A. Kricker, compares two knot--three-manifold pairs $(M_1,K_1)$ and $(M_2,K_2)$ ``of the above sort''. Let $\Sigma_p(M_i,K_i)$ denote the source's cyclic $p$-fold branched cover, with its specified deck action of $\Z/p\Z$. The conjecture is the equivalence between existence of an isomorphism
\[
\varphi:H_1(\Sigma_p(M_1,K_1);\Z)\longrightarrow
H_1(\Sigma_p(M_2,K_2);\Z)
\]
that is $\Z/p\Z$-equivariant and preserves linking, and equivalence of the pairs by a finite sequence of mod-$p$ loop moves.

With deck generators $\tau_i$ fixed, equivariance means $\varphi\tau_1=\tau_2\varphi$. For a closed oriented three-manifold the usual torsion linking form is $\ell_i:\operatorname{Tor}H_1\times\operatorname{Tor}H_1\to\Q/\Z$, and preservation means
\[
\ell_2(\varphi x,\varphi y)=\ell_1(x,y)
\]
for torsion classes $x,y$. The source requires an isomorphism on the full integral $H_1$, not just mod-$p$ homology or the torsion subgroup. If its covers are rational homology spheres, the whole $H_1$ is torsion, but that hypothesis is not reproduced and is not assumed silently.

The preceding class of pairs, admissible $p$, cover construction/meridian and orientation choices, and definition of a mod-$p$ loop move are absent. In a general manifold a cyclic branched cover requires appropriate homological data; it is not determined just by the notation without the source conventions. Equivariance also must not be weakened to allowing an unspecified automorphism of the deck group. Both directions of the stated iff remain part of the question, but no missing move or pair category is invented.
\subsection{Short English statement}
Do deck-equivariant, linking-preserving integral homology isomorphisms of the specified cyclic branched covers characterize equivalence by finitely many mod-$p$ loop moves? The pair category and move definition are missing.
\subsection{Sources}
\begin{itemize}
\item Local UnsolvedMath registry, permanent ID 10400043, exact source code catalogue entry, Conjecture 2.21; source attribution A. Kricker, no proposer/year in metadata. Registry status open, classification OPEN-TRIAGE; the registry title is truncated.
\item Ohtsuki, \emph{Problems on invariants of knots and 3-manifolds} (2002), Conjecture 2.21, PDF page 46, \url{https://msp.org/gtm/2002/04/gtm-2002-04-024s.pdf}. Supplied pdf-font-structure extraction accessed July 29, 2026; provenance distinguishes 231 bold numbered items. No PDF was accessed here.
\item The supplied summary identifies Kricker's mod-$p$ loop-move question but provides no precise later result or the preceding definitions.
\item Provenance flags current primary-literature status and redistribution terms NEEDS\_REVIEW.
\end{itemize}
\end{document}
